% Generated draft: requires mathematical and source review.
\documentclass[11pt]{article}
\usepackage[margin=25mm]{geometry}
\usepackage{fontspec}
\setmainfont{Latin Modern Roman}
\usepackage{amsmath,amssymb,mathtools,unicode-math}
\setmathfont{Latin Modern Math}
\usepackage{enumitem,needspace,xurl,hyperref}
\setlength{\parindent}{0pt}
\setlength{\parskip}{5pt}
\setlength{\emergencystretch}{3em}
\newcommand{\N}{\mathbb{N}}
\newcommand{\Z}{\mathbb{Z}}
\newcommand{\Q}{\mathbb{Q}}
\newcommand{\R}{\mathbb{R}}
\newcommand{\C}{\mathbb{C}}
\newcommand{\F}{\mathbb{F}}
\newcommand{\A}{\mathbb{A}}
\newcommand{\PP}{\mathbb{P}}
\newcommand{\E}{\mathbb{E}}
\newcommand{\eps}{\varepsilon}
\newcommand{\dd}{\,\mathrm{d}}
\begin{document}
\section*{UnsolvedMath ID: 5300035}
\textbf{Source problem code:} catalogue entry
\subsection{Definitions and mathematical statement}
A Siegel disk of a holomorphic dynamical system $f$ is a maximal rotation domain $U$ conformally equivalent to the unit disk, on which some iterate $f^p$, $p\geq1$, is conformally conjugate to an irrational rotation. Thus there exist $\alpha\in\mathbb R\setminus\mathbb Q$ and a conformal isomorphism $\phi:\mathbb D\to U$ such that
\[
f^p(\phi(z))=\phi(e^{2\pi i\alpha}z)\qquad(z\in\mathbb D).
\]
In the usual rational-map setting $U$ is a periodic Fatou component. A Jordan curve is the image of the circle under a continuous injective map, equivalently a subset homeomorphic to the circle.

Does there exist a Siegel disk $U$ whose boundary $\partial U$ is not a Jordan curve? The boundary is taken in the ambient dynamical surface (the Riemann sphere in the rational-map setting). The supplied item does not specify whether $f$ must be rational, polynomial, quadratic, or belong to a broader holomorphic class, so these potential versions are not identified.

The supplied research summary reports an affirmative answer attributed to Herman and refinements by Buff--Ch\'eritat, including asserted examples with non-locally-connected boundary. It supplies no precise map class, theorem citation, or complete arithmetic hypothesis. Its illustrative growth condition on continued-fraction denominators contains an unspecified exponent, $q_{n+1}\gg q_n^{?}$, and is not a usable theorem assumption. These resolution claims are recorded as supplied attributions rather than verified constructions.
\subsection{Short English statement}
Can the boundary of a Siegel disk fail to be a simple closed curve? The supplied summary reports such examples, but omits precise references and the class of maps involved.
\subsection{Sources}
\begin{itemize}
\item John Milnor, 1992 problem list, local problem P10, as collected in Stony Brook, \emph{Open Problems in Dynamical Systems}, \url{https://www.math.stonybrook.edu/open-problems-dynamical-systems}; accessed 2026-07-29, with source-TeX and PDF-text extraction recorded in the provenance.
\item Herman and Buff--Ch\'eritat, as the resolution attributions in the supplied research summary; no publication titles, dates, or theorem identifiers are supplied.
\item Local registry \texttt{data/unsolvedmath/problems.json}, permanent UnsolvedMath ID $5300035$, source code catalogue entry, \emph{Non-Jordan Siegel-disk boundary}. Dataset attribution: Ulam AI and the UnsolvedMath Contributors.
\end{itemize}
\end{document}
